\lysh{32753}{4}

\begin{alist}
\item Η σχέση μεταξύ του μήκους $y$ (σε cm) οστού του μηρού και του ύψους $x$ (σε cm) μιας
γυναίκας είναι: $y = 0,43x - 26$. Άρα για $y = 38,5$ έχουμε ισοδύναμα:
\begin{gather*}
38,5 = 0,43x - 26\text{, δηλαδή} \\
64,5 = 0,43x\text{, οπότε} \\
x = 150.
\end{gather*}
Επομένως το ύψος της γυναίκας είναι $150 \text{ cm}$.

\item Για να προέρχονται από τον ίδιο άνδρα το μηριαίο οστό και τα οστά χεριού, πρέπει το
μηριαίο οστό μήκους $y = 42,8 \text{ cm}$ να ανήκει στον άνδρα ύψους περίπου $x = 164 \text{ cm}$, που
ισχύει αφού: $0,45 \cdot 164 - 31 = 73,8 - 31 = 42,8$.

\item Θα εξετάσουμε αν υπάρχει τιμή της μεταβλητής $x$, ώστε: $0,43x - 26 = 0,45x - 31$.
Έχουμε ισοδύναμα:
\begin{gather*}
0,43x - 26 = 0,45x - 31\text{, οπότε} \\
5 = 0,02x\text{, δηλαδή} \\
x = 250.
\end{gather*}
Πρέπει δηλαδή ο άνδρας και η γυναίκα να έχουν ύψος $250 \text{ cm}$, που είναι φυσικώς αδύνατο.
Συνεπώς δεν μπορεί ένας άνδρας και μια γυναίκα ίδιου ύψους να έχουν μηριαίο οστό ίδιου
μήκους.
\end{alist}
